% ============================================================
% Theorem 4.1: SGD-Induced Sigma Suppression
% Status: FULL RECONSTRUCTION
% ============================================================

% --- Definitions and Assumptions ---
\begin{definition}[Schema Coherence ODE]
The schema coherence dynamics are governed by:
\[
\dot{\sigma}_A = \sigma_A \cdot
[\rho \cdot P_A \cdot \alpha_A \cdot (1 - \gamma_\sigma \cdot \sigma_A)
- \epsilon_\sigma \cdot \Omega_{SL}],
\]
as defined in Equation~\eqref{eq:sigma-ode} of the manuscript
and Equation~\eqref{eq:sigma-ode-recon} of the reconstruction.
\end{definition}

\begin{definition}[SGD Training Regime]
Standard SGD training minimizes a loss function $\mathcal{L}(\theta)$
on the training dataset, where $\theta$ denotes all network parameters.
No explicit $\sigma_A$-targeting term (e.g., a loss bonus for high
$\sigma_A$) is present. The shortcut-learning pressure
$\Omega_{SL}(d,t) \in [0, \Omega_{\max}]$ quantifies how much of the
task can be solved via surface-level pattern matching.
\end{definition}

\begin{assumption}[Shortcut Pressure Dominance in Early Training]
In the neighbourhood of $\sigma_A > 0$, SGD (optimizing purely for loss)
induces $\Omega_{SL} > 0$ whenever surface-level shortcuts exist.
Specifically:
\[
\epsilon_\sigma \cdot \Omega_{SL} > \rho \cdot P_A \cdot \alpha_A
\]
holds during early-to-mid training whenever the dataset contains
spurious correlations or surface-level patterns.
\end{assumption}

\begin{assumption}[Parameter Positivity]
All rate constants are non-negative: $\rho, \gamma_\sigma, \epsilon_\sigma > 0$,
and the state variables satisfy $P_A \in [0, P_{\max}]$,
$\alpha_A \in [0, 1]$, $\Omega_{SL} \in [0, \Omega_{\max}]$.
\end{assumption}

% --- Theorem Statement ---
\begin{theorem}[SGD-Induced $\sigma$-Suppression]
\label{thm:suppression}
Under standard SGD training (without explicit $\sigma_A$ targeting),
the schema coherence dynamics satisfy:
\[
\left.\frac{\partial \sigma_A(d,t)}{\partial t}\right|_{\text{SGD}} < 0
\]
in the neighbourhood of $\sigma_A > 0$ whenever the bracket in
\eqref{eq:sigma-ode-recon} is negative, i.e., when
\[
\epsilon_\sigma \cdot \Omega_{SL}(d,t) > \rho \cdot P_A(d,t) \cdot \alpha_A(d,t).
\]
In this regime, the $\sigma_A = 0$ equilibrium is locally stable and
attracts all nearby trajectories (the ``$\sigma$-trap''). When the
inequality reverses, $\sigma_A = 0$ becomes unstable and the system
bifurcates toward the schema-coherent branch (Proposition~4.1).
\end{theorem}

% --- Proof ---
\begin{proof}
The proof proceeds by sign analysis of \eqref{eq:sigma-ode-recon}
under the SGD regime.

\paragraph{Step 1: Sign of $\dot{\sigma}_A$ for $\sigma_A > 0$.}
From \eqref{eq:sigma-ode-recon}, for any $\sigma_A > 0$:
\[
\dot{\sigma}_A = \sigma_A \cdot B,
\]
where $B = \rho P_A \alpha_A (1 - \gamma_\sigma \sigma_A) - \epsilon_\sigma \Omega_{SL}$.
Since $\sigma_A > 0$ (strictly positive), the sign of $\dot{\sigma}_A$
is identical to the sign of $B$:
\[
\text{sgn}(\dot{\sigma}_A) = \text{sgn}(B).
\]

\paragraph{Step 2: SGD induces negative $B$ in the suppression regime.}
Under standard SGD training, the shortcut-learning pressure
$\Omega_{SL}$ is driven by the loss landscape. Because SGD greedily
minimizes training loss, it preferentially exploits surface-level
patterns when available, increasing $\Omega_{SL}$. In the
neighbourhood of $\sigma_A > 0$ (small but positive schema coherence),
the $(1 - \gamma_\sigma \sigma_A)$ factor is approximately 1, so:
\[
B \approx \rho P_A \alpha_A - \epsilon_\sigma \Omega_{SL}.
\]
When Assumption~1 holds ($\epsilon_\sigma \Omega_{SL} > \rho P_A \alpha_A$),
we have $B < 0$ and consequently $\dot{\sigma}_A < 0$.

\paragraph{Step 3: Local stability of $\sigma_A = 0$.}
Consider the linearization of \eqref{eq:sigma-ode-recon} at
$\sigma_A = 0$. The derivative of the right-hand side with respect
to $\sigma_A$ is:
\[
\frac{\partial \dot{\sigma}_A}{\partial \sigma_A}\bigg|_{\sigma_A = 0}
= [\rho P_A \alpha_A (1 - 2\gamma_\sigma \cdot 0) - \epsilon_\sigma \Omega_{SL}]
= \rho P_A \alpha_A - \epsilon_\sigma \Omega_{SL}.
\]
When $B < 0$ (the suppression regime), this eigenvalue is negative:
\[
\lambda = \rho P_A \alpha_A - \epsilon_\sigma \Omega_{SL} < 0.
\]
By the Hartman--Grobman theorem, the $\sigma_A = 0$ equilibrium is
locally exponentially stable. All trajectories starting sufficiently
close to $\sigma_A = 0$ converge to zero — this is the $\sigma$-trap.

\paragraph{Step 4: Global attraction within the basin.}
The ODE \eqref{eq:sigma-ode-recon} is one-dimensional in $\sigma_A$
(with other variables treated as parameters). The exact integral
(Proposition~3.2, Equation~\eqref{eq:exact-integral}) gives:
\[
\sigma_A(t) = \sigma_A(0) \cdot
\exp\!\left( \int_0^t B(s) \, ds \right).
\]
If $B(s) \leq -c < 0$ uniformly over $[0, t]$, then:
\[
\sigma_A(t) \leq \sigma_A(0) e^{-c t} \to 0 \text{ as } t \to \infty.
\]
Thus, whenever the bracket remains negative, $\sigma_A = 0$ is globally
attracting within the domain $\sigma_A \geq 0$.

\paragraph{Step 5: Bifurcation when the inequality reverses.}
When $\epsilon_\sigma \Omega_{SL} < \rho P_A \alpha_A$, the linearization
eigenvalue $\lambda$ becomes positive, making $\sigma_A = 0$ unstable.
The non-trivial equilibrium:
\[
\sigma_A^* = \frac{1}{\gamma_\sigma}\left(1 - \frac{1}{R_0}\right),
\quad R_0 = \frac{\rho P_A \alpha_A}{\epsilon_\sigma \Omega_{SL}},
\]
becomes positive and stable. This is the transcritical bifurcation
analyzed in Proposition~4.1 ($\sigma_{\text{critical}} = 1/\gamma_\sigma$).
$\square$
\end{proof}

% --- Boundary Cases ---
\begin{boundary-case}
\textbf{$\sigma_A(0) = 0$ exactly:}
By Proposition~3.2, $\dot{\sigma}_A = 0$ and $\sigma_A(t) \equiv 0$
(the invariant manifold). The $\sigma$-trap is permanent — no
bifurcation can occur because $\sigma_A$ is identically zero.

\textbf{$\sigma_A(0)$ very small but positive:}
The exact integral shows exponential decay to zero when $B < 0$.
The trajectory may take arbitrarily long to approach zero, but
asymptotically it converges. This is the $\sigma$-trap basin.

\textbf{$\Omega_{SL} = 0$ (no shortcut pressure):}
The inequality $\epsilon_\sigma \Omega_{SL} > \rho P_A \alpha_A$
cannot hold (unless $\rho P_A \alpha_A = 0$). SGD never suppresses
$\sigma_A$, and the system always evolves toward the
schema-coherent branch. This is the ideal training scenario.

\textbf{$\rho P_A \alpha_A = 0$ (no growth engine):}
The inequality holds trivially ($0 > 0$ is false, but $B = -\epsilon_\sigma \Omega_{SL} < 0$).
The $\sigma_A = 0$ equilibrium is stable regardless of $\Omega_{SL}$.
Without the growth engine ($\rho P_A \alpha_A$), schema coherence can
never increase.
\end{boundary-case}

% --- Circular Dependency Check ---
\begin{circularity-check}
\textbf{Dependencies on earlier results:} The proof uses
Proposition~3.2 (exact integral, forward invariance) and
Proposition~4.1 (bifurcation analysis). The dependency on
Proposition~4.1 is a forward reference for the bifurcation
mechanism, not for the suppression claim itself. The suppression
claim (Theorem~4.1) is established independently via sign analysis
of the bracket $B$.

\textbf{Forward dependencies:} Proposition~4.2 (depth threshold)
and Proposition~4.3 (hysteresis) depend on Theorem~4.1. These are
forward dependencies.

\textbf{Ordering concern:} The theorem statement references
Proposition~4.1 before Proposition~4.1 has been proved in the
reconstruction document. This is acceptable because the theorem
establishes the suppression regime independently; the reference
is to the bifurcation mechanism that operates when the inequality
reverses. In the reconstruction document, Proposition~4.1 appears
earlier in the ordering (Part~V vs. Part~IV). The cross-reference
is resolved.

\textbf{Self-circularity:} Assumption~1 (shortcut pressure dominance)
is specific to SGD training and is not formally derived from first
principles. It is a modelling assumption about the effect of SGD on
$\Omega_{SL}$. This is acceptable: the theorem proves that \emph{if}
SGD induces sufficient $\Omega_{SL}$, then $\sigma_A$ is suppressed.
The assumption is the bridge connecting the ODE analysis to the
SGD algorithm — a necessary modelling step.
\end{circularity-check}
